\section{Related Work}

\subsection{Sliders and steering}

Concept Sliders \citep{gandikota2024sliders} train a LoRA \citep{hu2022lora} so that its output matches the frozen model's prediction shifted along the difference between a positive-prompt and a negative-prompt prediction. SliderSpace \citep{gandikota2025sliderspace} drops the prompt pair, takes principal components of CLIP embeddings of many images of one concept, and trains one adapter per component. Measured Sliders \citep{chen2026measured} define image sliders by differentiable measurements and hold the other measurements still; we borrow their selectivity score.

In audio, FreeSliders \citep{ezra2025freesliders} apply the Concept Sliders target at inference time without training, across image, video, and audio models; most of their audio concepts are sound effects. TADA \citep{staniszewski2026tada} steers the cross-attention activations of audio diffusion transformers with vectors from prompt pairs, or from sparse-autoencoder features selected with prompt pairs. Facchiano et al.~\citep{facchiano2025patching} compute difference-in-means vectors for tempo and timbre; AnchorSteer \citep{chang2026anchorsteer} learns text-free concept modules on self-generated clips; Singh et al.~\citep{singh2026discovering} find and steer sparse features in autoregressive music models. All of these score a control by text--audio alignment and by perceptual distance to the unsteered output. TADA checks two concepts against a descriptor, and Wang \citep{wang2026counterfactual} shows with a key estimator and a beat tracker that apparent obedience to an instruction can be the model's prior. We are not aware of work that scores every control against a waveform property fixed in advance, compares it with plain signal processing, or reports what else moves.

\subsection{Where axes come from, and how much of them a model covers}

In image diffusion, principal components of a U-Net's bottleneck activations turn out to be semantic \citep{kwon2023hspace,haas2024directions}. Guinot et al.~\citep{guinot2026steering} learn sparse features of music embeddings to steer retrieval, and independent components of embeddings are known to read more cleanly than principal ones \citep{yamagiwa2023ica}. We found no work that takes axes from embeddings of real music and trains generator controls along them.

Diversity against real music is usually one number per model, such as recall and coverage in an embedding space \citep{novack2024presto}. Slendebroek and Metaxa \citep{slendebroek2026flattening} compare two commercial systems with human recordings on hand-built signal features and report homogenisation. Our comparison is per axis, on axes learned from real recordings, for open models that a slider can then be trained on. RIME \citep{schaffer2026rime} builds paired edit data by applying effects with known parameters; we use such effects the other way round, as a baseline a learned control has to beat.
